\section{Redundant (administrative) work}
\label{section:node-level:redundant-admin}


We have checked that

\begin{itemize}
  \item all of our source code is covered by parallelism;
  \item the workload is reasonably well-balanced;
  \item there is not much tasking overhead;
  \item the tasks are relatively long.
\end{itemize}


\noindent
If the code still does not scale particularly well, we may assume that something is flawed with the per-task/per-parallel region overhead:

\begin{hypothesis}
 Each task or parallel region, respectively, introduces some overhead which is constant.
 As the number of tasks increases, this overhead increases, too.
\end{hypothesis}


\noindent
Such an arrangement would materialise in an quasi-sequential program fraction $(1-f_{\text{intra}})$ in \eqref{equation:intra-node:performance-model}:
The code is parallelised, but this overhead scales with $p_{\text{threads}}$.
It therefore feeds into the sequential part of the scaling formula, i.e.~it is redundant work which we could also have done once outside of the parallel region.


\paragraph{Data collection}

If our hypothesis holds, then the measured $f_{\text{intra}}$ should decrease as we add more and more threads.
This might be an almost negligible effect. However, 
we should also see that the ratio of scientifically relevant calculations (cmp.~Section~\ref{section:preparation:labelling}) should decrease, while the remaining routines gain relative weight.
This means that we succeed in splitting up the scientifically relevant calculations among the threads, but the code increases in return the total amount of compute spent on administrative tasks.



\begin{instructions}{\vtune}
Run two standard hotspot analyses for two different thread counts.
Filter into the phase of interest and study the Bottom-up profile where you switch the Effective Time column into percentage.
The relative contribution of the most expensive (scientific) routines will be the lower the higher the thread count.
The effect typically becomes more pronounced if you narrow the filter down to one thread.   
\end{instructions}


\paragraph{Evaluation}



Potential optimisation

%Gibt es ein Amdahl law?
%   \item[$\boxplus$] Use Amdahl's law to characterise your code. If it fits to
%   the law, i.e.~there is an $f_{\text{intra}}$ such that
%   \eqref{equation:amdahl:speedup}.


\paragraph*{How it works}

\begin{itemize}[noitemsep]
  \item[$\boxplus$] Run a strong scaling test and plot the efficiency. Narrow down the number of setups of interest.
\end{itemize}


\paragraph*{How it does not work}

\begin{itemize}[noitemsep]
  \item[$\boxminus$] Start with a full-node run right from the start and hope that you spot interesting trends in the data.
\end{itemize}
